% Standalone editor preview. When included in the paper, use its document class and bibliography.
\ifdefined\subsection
\else
  \documentclass[10pt]{article}
  \usepackage[margin=1in]{geometry}
  \renewcommand{\thesection}{\Roman{section}}
  \renewcommand{\thesubsection}{\Alph{subsection}}
  \def\eclipseFinishPreview{\bibliographystyle{unsrt}\bibliography{SCITECH_IA_Background,SCITECH_IIA_Technologies,SCITECH_IIA_Axe5_Fondements,SCITECH_IIA_Literature_Review}\end{document}}
  \begin{document}
  \section{Introduction}
\fi

\subsection{Background}

NASA's plans for sustained lunar exploration envisage longer surface stays and repeated access to sites of scientific interest~\cite{nasa_moonbase_faq_2026}. Returning to a site would allow crews to pursue questions raised by earlier samples, extend field investigations into surrounding terrain, and assess the potential uses of local resources. Successive missions would also provide opportunities to test technologies under lunar conditions and gain experience maintaining equipment far from Earth in preparation for human exploration of Mars~\cite{merancy_moon_to_mars_2025}. Infrastructure that remains available between crewed visits could support these objectives as the frequency and duration of missions increase~\cite{nasa_moonbase_faq_2026,merancy_moon_to_mars_2025}.

Reusing infrastructure across missions changes what engineers must consider when designing each addition. A power system deployed for an initial habitat might later serve a resource-processing facility or support surface mobility. Whether it can meet these demands depends on where the new assets are placed, how much power is available, and which operations must run concurrently. As the base develops, systems with distinct functions and management responsibilities increasingly rely on shared services~\cite{nasa_moonbase_guide_2026,darpa_luna10_results_2024}. Their ability to accomplish a mission together makes the base a system of systems and calls for evaluating their interactions over the course of operations.

Repeated access to the surface brings these interactions into focus at a lunar spaceport, a logistical and operational hub considered in the LUNAverse initiative~\cite{mueller_lunaverse_preprint_2026}. A site suitable for landing may offer limited access to illuminated terrain or exploitable resources~\cite{nasa_moonbase_faq_2026}. The relative placement of landing areas, power assets, and processing facilities therefore affects traverse routes, energy demand, and the operations that can be sustained~\cite{mueller_lunaverse_preprint_2026,speyerer_traverse_2016}. Studying how these systems would operate together can inform early decisions about which capabilities to deploy first and which interfaces they must support~\cite{merancy_moon_to_mars_2025}. To evaluate a common configuration, specialists in terrain, mobility, power, and processing must relate their models to the same sites, assets, and operating schedules. Where separate organizations develop those systems, establishing compatible assumptions requires collaboration across organizational as well as disciplinary boundaries~\cite{darpa_luna10_results_2024,mueller_lunaverse_preprint_2026}.

\subsection{Motivation}

The difficulty of this collaboration stems in part from how engineering knowledge is distributed across specialized organizations in the space sector~\cite{oecd_space_2022}. Teams develop their models using different tools, data formats, and internal standards~\cite{ciampa_agile_2020}. Even within a single organization, multiple versions of the same system model may coexist, reflecting successive design iterations or different assumptions. Integrating these contributions requires establishing which models, interfaces, and assumptions are compatible. This challenge extends beyond technical interoperability: organizations may need to exchange system specifications and performance data while protecting proprietary models, intellectual property, and commercially sensitive information~\cite{aisenberg_rights_2014}. Building a common architecture therefore requires coordinating heterogeneous engineering contributions while respecting the constraints on what each participant can share.

Beyond these integration challenges, the combination of design alternatives across the many interdependent systems of a lunar infrastructure creates a potentially vast architectural design space. Different technologies, capacities, quantities, layouts, and operational strategies may satisfy the same mission requirements, while their interactions constrain which configurations are feasible~\cite{boyer_delos_2025}. This diversity motivates a modular approach in which systems and their behaviors can be selected, modified, and recombined without rebuilding the entire engineering process for each configuration. A shared engineering environment would help designers assemble compatible alternatives from the information partners can disclose, compare their performance, and select an architecture suited to the mission requirements.

Evaluation itself presents a distinct challenge. Quantitative performance measures are essential for assessing mission outcomes, but they do not always reveal how those outcomes emerge from interactions between systems. Understanding the progression of a scenario may require examining when systems operate, how resources are exchanged, and where operational constraints or bottlenecks arise. Because these interactions also depend on the spatial arrangement of assets and the timing of their activities, spatial and temporal visualization can complement quantitative analysis by allowing engineers to examine simulated operations within their operational context.

Digital engineering provides an integrated, model-based approach to managing these challenges. By maintaining a consistent digital thread connecting requirements, system definitions, architectural configurations, simulations, and evaluation results, digital engineering ecosystems can support multidisciplinary design and iterative analysis~\cite{dod_digital_engineering_2023}. Developed within the LUNAverse initiative~\cite{mueller_lunaverse_preprint_2026}, ECLIPSE applies this approach to lunar surface infrastructure architecting. An authoritative Systems Modeling Language (SysML) model relates mission requirements to system definitions and interfaces, supporting the selection and configuration of modules for different scenarios. These scenarios are then simulated to assess their operational performance, while three-dimensional visualization in NVIDIA Omniverse enables their spatial and temporal evolution to be examined within a representation of the lunar terrain. This integrated workflow is demonstrated through a Lunar Spaceport-1 in situ resource utilization scenario, connecting system configuration, operational simulation, quantitative evaluation, and synchronized spatial playback.

